% ==================================================================
% SECTION 3: METHODS
% ==================================================================
%
% This section describes the experimental methodology, including:
% - Datasets and data splits
% - Attack scenarios
% - Defense baselines
% - Evaluation metrics
% - Experimental setup
%
% Author: Luminous Dynamics
% Date: November 8, 2025
% Status: Template ready for results integration
% ==================================================================

\section{Experimental Methodology}
\label{sec:methods}

We evaluate our proposed defenses against a comprehensive suite of Byzantine attacks across multiple datasets and data distribution scenarios. This section details our experimental setup, attack models, defense baselines, and evaluation metrics.

% ------------------------------------------------------------------
% 3.1 Datasets and Data Splits
% ------------------------------------------------------------------

\subsection{Datasets}
\label{sec:datasets}

We evaluate on three federated learning datasets representing different complexity levels:

\textbf{EMNIST Balanced~\cite{cohen2017emnist}}. Extended MNIST with 47 balanced classes (10 digits + 37 letters), comprising 112,800 training images and 18,800 test images. Each grayscale image is 28×28 pixels. EMNIST balances class distribution across digits and uppercase/lowercase letters (excluding 15 confusing characters), making it ideal for evaluating detection under balanced label distributions.

\textbf{FEMNIST~\cite{caldas2018leaf}}. Federated Extended MNIST with 62 classes (10 digits + 26 uppercase + 26 lowercase letters) from 3,550 writers, comprising 671,585 training images. FEMNIST preserves natural data heterogeneity where each writer (client) has unique handwriting style and non-uniform label distribution, making it the most realistic federated dataset.

\textbf{MNIST~\cite{lecun1998mnist}}. Classic 10-class digit recognition with 60,000 training images and 10,000 test images. Used as a lightweight baseline for rapid prototyping and ablation studies.

% ------------------------------------------------------------------
% 3.2 Data Distribution Strategies
% ------------------------------------------------------------------

\subsection{Data Distribution Strategies}
\label{sec:data_distribution}

To simulate realistic heterogeneous federated settings, we consider two data partitioning strategies:

\textbf{IID (Independent and Identically Distributed)}. Training data is uniformly shuffled and partitioned into $N$ equal-sized subsets, one per client. Each client has the same label distribution as the global dataset. This represents the idealized case where clients have homogeneous data.

\textbf{Non-IID (Dirichlet Distribution)}. We use Dirichlet distribution with concentration parameter $\alpha$ to control label skew~\cite{hsu2019measuring}. For each client $i$ and class $k$, we sample $p_{i,k} \sim \text{Dir}(\alpha)$ and assign data points proportionally. Lower $\alpha$ values create higher skew:
\begin{itemize}[leftmargin=*,noitemsep]
    \item $\alpha = 0.1$: Extreme skew (1-3 classes per client)
    \item $\alpha = 0.3$: High skew (3-5 classes per client) \textbf{[Default]}
    \item $\alpha = 1.0$: Moderate skew (most classes present)
\end{itemize}

Unless stated otherwise, we use $\alpha = 0.3$ for non-IID experiments, representing a realistic federated scenario where clients have specialized but overlapping data distributions.

% ------------------------------------------------------------------
% 3.3 Attack Scenarios
% ------------------------------------------------------------------

\subsection{Byzantine Attack Models}
\label{sec:attacks}

We evaluate against 8 representative Byzantine attacks spanning three categories:

\subsubsection{Untargeted Attacks (Detection Focus)}

These attacks aim to degrade global model performance without targeting specific outcomes, making them ideal for evaluating detection effectiveness.

\textbf{Sign Flip}~\cite{fung2018mitigating}. Byzantine clients submit $-g_i$ where $g_i$ is the honest gradient. This simple attack is highly effective against naive averaging but easily detectable.

\textbf{Scaling Attack (×100)}. Byzantine clients submit $100 \times g_i$, attempting to dominate aggregation through magnitude. Tests defenses' robustness to gradient norm manipulation.

\textbf{Random Noise}. Byzantine clients submit $\mathcal{N}(0, \sigma^2 I)$ where $\sigma = ||\bar{g}||$ matches the average honest gradient norm. Tests detection of statistically unusual but plausible gradients.

\textbf{Noise Masking}. Byzantine clients add scaled noise to honest gradients: $g_i + \eta \mathcal{N}(0, I)$ where $\eta$ is chosen to maintain plausible gradient norms. More sophisticated variant of random noise.

\textbf{Targeted Neuron Poisoning}. Byzantine clients zero out gradients for 90\% of parameters and amplify the remaining 10\% by factor of 10. Tests detection of sparse, high-magnitude perturbations.

\subsubsection{Collusion Attacks (Coordination)}

\textbf{Collusion Attack}. Byzantine clients coordinate to submit identical malicious gradients, testing defenses against Sybil behavior without explicit identity duplication.

\textbf{Sybil + Sign Flip}. Combines identity forgery with sign-flip attack. Multiple Byzantine identities submit $-g_i$, testing Sybil resistance of aggregation mechanisms.

\subsubsection{Adaptive Attacks (Evasion)}

\textbf{Sleeper Agent}. Byzantine clients behave honestly for $W$ warmup rounds before activating attack (sign flip), testing defenses' temporal consistency tracking and warm-up mechanisms.

% ------------------------------------------------------------------
% 3.4 Defense Baselines
% ------------------------------------------------------------------

\subsection{Defense Baselines}
\label{sec:defenses}

We compare our proposed defenses (CBF, PoGQ-v4.1) against 6 state-of-the-art Byzantine-robust aggregation methods:

\textbf{FedAvg}~\cite{mcmahan2017communication}. Undefended baseline averaging all client updates: $g_{\text{agg}} = \frac{1}{N}\sum_{i=1}^{N} g_i$. Highly vulnerable to all attacks but serves as lower bound for detection performance.

\textbf{Coordinate-Wise Median (Coord-Median)}~\cite{yin2018byzantine}. Computes median independently for each parameter: $g_{\text{agg}}[j] = \text{median}(g_1[j], \ldots, g_N[j])$. Robust to $\lfloor N/2 \rfloor$ Byzantine clients but provides no detection mechanism.

\textbf{Coordinate-Wise Median with Safety Guards (CM-Safe)}. Our hardened variant of Coord-Median with 4 safety guards:
\begin{itemize}[leftmargin=*,noitemsep]
    \item \textit{Min-Clients Guard}: Rejects aggregation if fewer than 50\% honest clients
    \item \textit{Norm-Clamp Guard}: Clamps gradient norm to safe range
    \item \textit{Direction-Check Guard}: Verifies aggregated direction vs.\ honest median
    \item \textit{Reject-All Fallback}: Returns zero update if all guards fail
\end{itemize}

\textbf{Robust Federated Aggregation (RFA)}~\cite{pillutla2019robust}. Uses geometric median with smoothed Weiszfeld algorithm. Provably robust to $\lfloor N/2 \rfloor$ Byzantine clients with $O(N \log N)$ computational overhead.

\textbf{FLTrust}~\cite{cao2020fltrust}. Server maintains root dataset and uses cosine similarity between server gradient and client gradients for trust scoring. Requires clean server-side data.

\textbf{BOBA}~\cite{rodriguez2022boba}. Byzantine-robust aggregation with probabilistic analysis. Combines coordinate-wise filtering with randomized client selection.

\textbf{Conformal Byzantine Filter (CBF)}. Our Phase-1 detector combining PCA-based dimensionality reduction with conformal prediction for principled FPR control. Implements:
\begin{itemize}[leftmargin=*,noitemsep]
    \item PCA projection to top-$k$ components (default $k=10$)
    \item Conformal score: $s_i = ||g_i - \bar{g}_{\text{honest}}||_2$ in PCA space
    \item Rejection threshold: $q_{\alpha}$ from calibration set, guaranteeing FPR $\leq \alpha$ up to quantile estimation error (under exchangeability assumption; non-IID data may exhibit calibration drift) (default $\alpha = 0.10$)
\end{itemize}

\textbf{PoGQ-v4.1 (Ours)}. Our flagship multi-component detector integrating:
\begin{enumerate}[leftmargin=*,noitemsep]
    \item \textit{PCA Extractor}: Top-10 principal components (90\% variance)
    \item \textit{Conformal Thresholding}: $\alpha = 0.10$ with 2\% margin
    \item \textit{Mondrian Conformal}: Per-bucket thresholds for 3 profiles (label, norm, direction)
    \item \textit{Adaptive Hybrid}: Switches between Mondrian and global based on bucket size
    \item \textit{EMA Smoothing}: $\beta = 0.85$ exponential moving average
    \item \textit{Hysteresis}: $(k=2, m=3)$ prevents quota flapping
\end{enumerate}

All defenses operate in the honest-majority setting where Byzantine clients constitute $\leq 40\%$ of total clients (unless otherwise noted).

% ------------------------------------------------------------------
% 3.5 Evaluation Metrics
% ------------------------------------------------------------------

\subsection{Evaluation Metrics}
\label{sec:metrics}

We evaluate defense performance using both \textit{detection metrics} and \textit{model utility metrics}.

\subsubsection{Detection Metrics}

\textbf{Area Under ROC Curve (AUROC)}. Primary metric measuring discrimination between Byzantine and honest clients. AUROC = 1.0 indicates perfect detection; AUROC = 0.5 indicates random guessing.

\textbf{True Positive Rate at Fixed FPR (TPR@FPR)}. We report TPR at target FPR = 0.10, representing the fraction of Byzantine clients detected while maintaining $\leq 10\%$ false positive rate on honest clients.

\textbf{False Positive Rate at Fixed TPR (FPR@TPR)}. Conversely, we report FPR at target TPR = 0.90, measuring honest client rejection rate when detecting 90\% of attacks.

\textbf{Time-to-Detection (TTD)}. Median number of rounds until a Byzantine client is first flagged. Lower is better. We define TTD = $-1$ if a Byzantine client is never detected within the experiment duration.

\textbf{Per-Bucket False Positive Rate}. For Mondrian conformal detectors, we verify that FPR $\leq \alpha$ holds \textit{for each bucket independently} up to quantile estimation error (under exchangeability within buckets), ensuring calibration holds across data heterogeneity.

\subsubsection{Model Utility Metrics}

\textbf{Test Accuracy}. Classification accuracy on held-out test set after $T$ rounds of federated training.

\textbf{Convergence Speed}. Number of rounds to reach target accuracy threshold (e.g., 90\% for MNIST, 80\% for EMNIST).

\textbf{Accuracy Degradation under Attack}. $\Delta_{\text{acc}} = \text{Acc}_{\text{clean}} - \text{Acc}_{\text{attack}}$ measures attack impact on final model quality.

\subsubsection{Statistical Significance}

All experiments are repeated with 2 random seeds. We report:
\begin{itemize}[leftmargin=*,noitemsep]
    \item \textbf{Mean $\pm$ Std Dev} for AUROC, TPR, FPR
    \item \textbf{Bootstrap 95\% Confidence Intervals} (1000 samples) for ROC curves
    \item \textbf{Wilcoxon Signed-Rank Test} ($p < 0.05$) for pairwise defense comparisons
\end{itemize}

% ------------------------------------------------------------------
% 3.6 Experimental Setup
% ------------------------------------------------------------------

\subsection{Experimental Setup}
\label{sec:experimental_setup}

\subsubsection{Federated Learning Configuration}

\textbf{Model Architecture}. We use a simple CNN architecture for all datasets:
\begin{itemize}[leftmargin=*,noitemsep]
    \item Conv2D(1, 32, 3×3, ReLU) $\rightarrow$ MaxPool(2×2)
    \item Conv2D(32, 64, 3×3, ReLU) $\rightarrow$ MaxPool(2×2)
    \item Flatten $\rightarrow$ Dense(128, ReLU) $\rightarrow$ Dense($C$, Softmax)
\end{itemize}
where $C$ is the number of classes (10 for MNIST, 47 for EMNIST, 62 for FEMNIST).

\textbf{Training Hyperparameters}.
\begin{itemize}[leftmargin=*,noitemsep]
    \item Local optimizer: SGD with learning rate $\eta = 0.01$, momentum $= 0.9$
    \item Local epochs per round: $E = 1$
    \item Batch size: 64
    \item Global rounds: $T = 20$
    \item Number of clients: $N = 100$
    \item Clients per round: 20 (20\% participation)
    \item Byzantine ratio: 20\% (4 out of 20 selected clients per round)
\end{itemize}

\textbf{Calibration Set}. For conformal prediction, we reserve 20\% of honest training data as a calibration set, ensuring proper coverage guarantees. Calibration is performed once before training begins.

\subsubsection{Computational Environment}

All experiments are conducted on:
\begin{itemize}[leftmargin=*,noitemsep]
    \item CPU: Intel Xeon (details TBD based on actual hardware)
    \item GPU: NVIDIA RTX 3090 (24GB VRAM)
    \item Framework: PyTorch 2.1.0, Python 3.11
    \item OS: NixOS 25.11 with reproducible environment
\end{itemize}

\textbf{Reproducibility}. We provide:
\begin{itemize}[leftmargin=*,noitemsep]
    \item Full source code at \url{https://github.com/luminous-dynamics/zero-trustml} (to be released)
    \item Nix flake for exact environment reproduction
    \item Experiment configs (YAML) for all reported results
    \item Random seed manifests for each experiment run
    \item SHA-256 hashes of all artifacts
\end{itemize}

\subsubsection{Experimental Matrix}

We evaluate two experimental slices:

\textbf{Sanity Slice} (256 experiments). Rapid validation across core scenarios:
\begin{itemize}[leftmargin=*,noitemsep]
    \item Datasets: EMNIST, MNIST
    \item Attacks: All 8 attacks
    \item Defenses: All 8 defenses
    \item Seeds: 2 random seeds
    \item Total: $2 \times 8 \times 8 \times 2 = 256$ experiments
\end{itemize}

\textbf{Full FEMNIST Matrix} (864 experiments). Comprehensive evaluation on realistic federated dataset:
\begin{itemize}[leftmargin=*,noitemsep]
    \item Dataset: FEMNIST (non-IID, $\alpha = 0.3$)
    \item Attacks: All 8 attacks
    \item Defenses: All 8 defenses
    \item Byzantine ratios: 20\%, 30\%, 40\%
    \item Seeds: 3 random seeds (increased for robustness)
    \item Total: $1 \times 8 \times 8 \times 3 \times 3 = 576$ base experiments
    \item Ablations: $\alpha \in \{0.1, 1.0\}$, $\beta \in \{0.85, 0.95\}$, $k/m$ variants
    \item Total with ablations: 864 experiments
\end{itemize}

\subsubsection{Proof Backend Architecture}

We implement a dual-backend STARK proof system with RISC Zero zkVM~\cite{risc0} as the primary backend and Winterfell AIR~\cite{winterfell} as an optional specialized backend. This architecture balances development velocity, soundness guarantees, and performance.

\textbf{RISC Zero Backend (Primary).} The primary implementation uses RISC Zero zkVM (v3.0), which provides a general-purpose RISC-V zero-knowledge virtual machine. The PoGQ decision logic is implemented in Rust and compiled to RISC-V, with proofs generated via the in-process prover. RISC Zero provides 127-bit conjectured security~\cite{risc0-security} with FRI blowup factor 4 (S128 profile, default) or factor 2 (S192 profile, 191-bit security). Proof generation takes 35.8 seconds per decision on consumer hardware (Intel Core i9-8950HK, 2.9 GHz, 6-core mobile CPU), generating 216 KB proofs with <100ms verification time. While 30,000$\times$ slower than domain-specific Winterfell AIR (1.5ms), RISC Zero's general-purpose RISC-V execution model provides cross-validation: identical guest code on both backends proves semantic correctness of our custom AIR implementation.

\textbf{Provenance Integration.} To ensure tamper-evident configuration commitment, we compute a Blake3~\cite{blake3} cryptographic hash over build metadata: Rust compiler version, Git commit hash, build timestamp, security profile ID (S128=128, S192=192), and AIR schema revision (monotonic version for constraint system changes). This 256-bit hash is committed in the proof's public inputs and verified via fail-fast assertions in the guest program. If profile\_id $\notin \{128, 192\}$ or air\_rev $\neq$ EXPECTED\_REV, the prover panics with a descriptive error, preventing proofs from mismatched configurations. Provenance is enabled via \texttt{VSV\_PROVENANCE\_MODE=strict} environment variable, allowing graceful degradation when not needed (e.g., local testing).

\textbf{Winterfell Backend (Optional).} The Winterfell AIR backend (v0.13.1) uses domain-specific constraints for 7-10$\times$ faster proving but encountered a fundamental limitation: AIR trace length must be determined at compile-time (static constraints), while PoGQ's variable-length gradient vectors require dynamic trace sizing. We explored selector column patterns and padding strategies but could not resolve this without sacrificing soundness. Winterfell remains 95\% complete and serves as a performance reference, but RISC Zero's flexibility makes it the production choice.

\textbf{Security Analysis.} RISC Zero's 127-bit security (S128 profile) exceeds the 128-bit standard by a narrow margin due to FRI security analysis~\cite{fri-security}. While this is slightly below the canonical 128-bit threshold, it provides sufficient security for FL audit scenarios where: (1) proofs are verified immediately (no long-term storage), (2) attack surface is limited to FL participants (not public), (3) worst-case failure is audit log tampering (not key compromise). For high-security deployments, S192 profile provides 191-bit security at 2$\times$ proving cost.

Cryptographic receipts attest \textit{decision correctness} under PoGQ logic---that quarantine state transitions follow specified rules---but do not attest raw gradient provenance or client identity (which require orthogonal mechanisms like secure enclaves or digital signatures).

% ------------------------------------------------------------------
% 3.7 Artifact Generation
% ------------------------------------------------------------------

\subsection{Artifact Generation and Provenance}
\label{sec:artifacts}

Each experiment run generates 6 JSON artifacts for transparency and reproducibility:

\begin{enumerate}[leftmargin=*,noitemsep]
    \item \texttt{detection\_metrics.json}: AUROC, TPR, FPR, ROC curve, time-to-detection
    \item \texttt{per\_bucket\_fpr.json}: Per-bucket FPR for Mondrian conformal (CI gate)
    \item \texttt{ema\_vs\_raw.json}: EMA flap counts and stability metrics (CI gate)
    \item \texttt{coord\_median\_diagnostics.json}: Guard activation counts (CI gate)
    \item \texttt{bootstrap\_ci.json}: Bootstrap confidence intervals (1000 samples)
    \item \texttt{model\_metrics.json}: Test accuracy, loss, convergence statistics
\end{enumerate}

Additionally, we generate:
\begin{itemize}[leftmargin=*,noitemsep]
    \item \texttt{MANIFEST.json}: Git commit, config hash, artifact SHA-256 hashes, optional Ed25519 signature
    \item \texttt{CALIBRATION.json}: PCA params, conformal quantiles, Mondrian profiles (with SHA-256 hash)
\end{itemize}

All artifacts are cryptographically signed (Ed25519) and version-controlled for audit compliance.

% ------------------------------------------------------------------
% 3.8 CI Gates (Hard Guarantees)
% ------------------------------------------------------------------

\subsection{Continuous Integration Gates}
\label{sec:ci_gates}

To ensure experimental rigor, we enforce 3 hard CI gates that must pass before accepting results:

\textbf{Gate 1: Conformal FPR Compliance}. For all Mondrian buckets, empirical FPR $\leq 0.12$ (target $\alpha = 0.10$ + 2\% margin). Violations indicate miscalibration.

\textbf{Gate 2: EMA Stability}. Hysteresis flap count $\leq 5$ per experiment (20 rounds). High flap counts indicate quota instability.

\textbf{Gate 3: CM-Safe Guard Coverage}. For Coord-Median-Safe experiments, $\geq 1$ guard must activate during Byzantine rounds. Zero activations indicate guard failure.

Experiments failing any gate are flagged for manual review and excluded from final results if unresolved.

% ------------------------------------------------------------------
% End of Methods Section
% ------------------------------------------------------------------

% ==================================================================
% NOTES FOR FUTURE COMPLETION
% ==================================================================
%
% TODO (after sanity slice completes):
% [ ] Fill in actual hardware specs (CPU model, RAM)
% [ ] Add example ROC curve figure reference
% [ ] Add example TTD histogram figure reference
% [ ] Include statistical significance results table
% [ ] Cite final GitHub repository URL
%
% TODO (after full FEMNIST matrix completes):
% [ ] Update ablation study details with actual parameters tested
% [ ] Add cross-dataset comparison subsection if results warrant
%
% ==================================================================
